% LSPV-4 report 384 2026-11-04 04:45 PST
\documentclass[11pt]{article}
\usepackage{amsmath,amssymb,amsthm,hyperref}
\newtheorem{theorem}{Theorem}
\title{LSPV-4: Resolution DAG Normalization and Read-Only Rescan Verification}
\author{PNP Time-Space Resource Auditor}
\date{November 4, 2026, 04:45 PST}
\begin{document}\maketitle
\noindent\textbf{PROVED/RESTRICTED}, not an unrestricted SAT complexity result. Fixed $R_{\rm mix341}$ unchanged. B1, B2, B3 and P versus NP OPEN.
\section*{Established hard-family anchor}
Haken (1985): for standard explicit unary pigeonhole $\mathrm{PHP}_{h+1}^h$ with $h(h+1)$ variables, $(h+1)+h\binom{h+1}{2}$ clauses and $\Theta(h^3\log(h+2))$ signed binary-index bits, there exist $c>0,h_0$ such that for all $h\ge h_0$ EVERY ordinary shared-DAG Resolution refutation has $2^{ch}$ or more lines. This does not concern Extended Resolution, GF(2), combined $R_{\rm mix341}$, or arbitrary SAT machines. Cook (1976) supplies polynomial Extended Resolution PHP proofs; odd Tseitin has affine GF(2) cancellations.
\section*{Exact restricted verifier theorem}
Fix strict ordinary Resolution: every input clause is a sorted, non-tautological signed list; each derived record names two earlier clauses, an opposing signed pivot, and the EXACT non-tautological resolvent. Each derived node costs one. Every source reference must authenticate to the explicit original CNF; no hidden semantics or auxiliaries. Let $N$ source bits, $q$ derived nodes, $P$ encoded proof bits and $M=N+P+O(1)$.
\begin{theorem}
Every $q$-node strict Resolution DAG refutation has a canonical equivalent with the SAME $q$ derived nodes and $P=O((N+q+2)^3)$ bits. One fixed deterministic indexed read-only-input Turing machine verifies any candidate of combined length $M$ with $O((M+2)^6)$ time and $O(\log(M+2))$ writable bits, rejecting malformed and unsound evidence. This is a provided-proof verifier, not a proof constructor.
\end{theorem}
\begin{proof}
All resolvents use original variables and at most $N$ signed literals, each with $O(\log(N+2))$ bits. Remove unused source-reference records and topologically reindex shared nodes. Encode each of at most $N+q$ records in $O(N\log N+\log(N+q))$ bits, giving the conservative cubic estimate. The indexed verifier keeps only $O(\log M)$-bit positions, counters, node IDs and literal IDs. For every record, it repeatedly rescans immutable source/proof input to check grammar, source authenticity, earlier references, pivot signs, exact set equality of the conclusion with the two parents minus resolved pivot, and absence of duplicates and complement pairs. Nested scans have constant nesting depth and fit the loose $O(M^6)$ native-step bound including index/head operations. An induction over the proof list proves that accepted conclusions follow from the original source, so SAT sources cannot have an accepted final empty clause.
\end{proof}
\section*{Falsification and full resource firewall}
\begin{enumerate}
\item[] F1. Claim that $q$ DAG nodes force $\Omega(q)$ simultaneous writable checker cells: FALSIFIED by indexed rereads, though not on every sequential-access model.
\item[] F2. Claim that node count bounds proof-validation cost without a grammar: FALSIFIED for arbitrary opaque payloads; normalized encoding exists here only for strict Resolution.
\item[] F3. Claim that $O(\log M)$ checker space gives $O(\log N)$ SAT decision space: FALSIFIED as a transfer. The proof is separately supplied; a solver must pay proof construction, source access, serialization and all-phase peak writable space.
\item[] F4. Claim that Haken gives exponential arbitrary-SAT machine space: FALSE proof-system transfer, contradicted as a general inference by Cook Extended Resolution and GF(2).
\end{enumerate}
Exponentially large proofs are NOT excluded: their external input $M$ and checking time can be exponential. Repeated read-only access substitutes time for writable storage; the input is NOT a free solver-generated proof. If materialized by the solver, its output alone costs $\Omega(P)$ bit writes. Discovery of a shortest derivation is omitted from $W_R$ by design, not from any actual upper-bound construction.
\section*{Reproducibility, peer review and next gap}
Python producer: 45 Horn-chain UNSAT proofs of cost $k+1$ for $k=0,\ldots,39,48,64,96,128,192$; SAT deletion controls; 45 forged pivots, 45 future references, 10 malformed cases rejected; zero-inference empty source accepted. Signed $K_4$ odd/even parity controls 0/8 models. 220 checks and raw SHA-256 \texttt{c02cdb282907b9a756b20d6b8ed3121ac7f6d3a58b94485831b310078f5022f7}. Independent Node 776 checks over 257 chains. Python host RSS 96532 KiB is NOT claimed machine space.
Reviewed verified 2026-11-02 1204 handoff (blob c3889a3e), current issue \#3, AGENTS.md, CODEX.md, pnp-research skill, LSPV-4/360 and /383, LSPV-3/392, LSPV-5/317 and /335. Grok actual \texttt{GROKPROOF/2026-11-04\_1217\_PROGRESS.md} (cae0560d) leaves the fully charged universal transformation open; its claimed cut complexity is restricted. Claude actual \texttt{research/agents/claude/2026-10-28\_1521\_CLAUDE\_ONETAPE-0001.tex} (9a5b05d6) concerns sequential-input lower bounds, not indexed FLMV. Project and OpenAI GitHub searches found no applicable full Rmix341 proof-normalizer. Relevant external verifiers: drat-trim, lean-sat-core, short-proof. No claim of research novelty for a standard repeated-scan construction.
The exact next obligation is full authoritative typed LOCAL-CONVERT/DEF/GF(2) grammar, and a representation/verification bound for every fixed mixed rule. SAT in P implying polynomial $\overline W_{R_{\rm mix341}}$ is OPEN.
\begin{thebibliography}{9}
\bibitem{Haken} A. Haken, The Intractability of Resolution, TCS 39 (1985), 297--308.
\bibitem{Cook} S. Cook, A Short Proof of the Pigeon Hole Principle Using Extended Resolution, SIGACT News 8(4) (1976), 28--32.
\bibitem{CR} S. Cook and R. Reckhow, The Relative Efficiency of Propositional Proof Systems, JSL 44 (1979), 36--50.
\bibitem{FLMV} L. Fortnow, R. Lipton, D. van Melkebeek, A. Viglas, Time-space lower bounds for satisfiability, JACM 52 (2005), 835--865.
\end{thebibliography}
\end{document}
